% =========================================================================
% APPENDICE A — ETICA, CONFINI E LIMITI
% =========================================================================

\chapter{Etica, confini e limiti}
\label{app:ethics}

Questa appendice consolida, in un unico luogo, i confini metodologici entro cui vanno lette le conclusioni di questa tesi. È organizzata attorno alle tre fasi dell'arco rilevare--tradurre--valutare. Ogni voce enuncia un confine, la ragione per cui è stato scelto e --- dove rilevante --- la direzione in cui un lavoro futuro potrebbe estendersi oltre esso. L'intento è la trasparenza: rendere leggibili i limiti del lavoro anziché disperderli nel testo principale.

% -------------------------------------------------------------------------
\section{Rilevare --- fase analitica}
\label{app:limits-detect}

\subsection*{Rilevamento di pattern contro predizione}
La fase analitica recupera struttura latente e caratterizza le anomalie; non fa previsioni. L'analisi dei segnali di preallarme restituisce un risultato nullo una volta corretta la dimensione campionaria effettiva per l'autocorrelazione delle finestre sovrapposte: un trend monotòno apparentemente forte si rivela un artefatto di pseudo-replicazione. Questo è riportato come contributo all'onestà metodologica, non come evidenza a favore o contro una transizione imminente. Nessuna affermazione sullo stato futuro dell'AMOC è fatta o implicata dalla fase analitica; il livello di previsione rinviato (O1.4) è solo specifica.

\subsection*{Osservazione contro ricostruzione}
L'analisi attinge a due fonti di diverso status epistemico. L'array RAPID è osservazione diretta; l'ensemble Copernicus GREP è una rianalisi --- output di modello vincolato dall'assimilazione dei dati. Le affermazioni che poggiano sulla rianalisi portano la dipendenza aggiuntiva da un modello, e sono formulate di conseguenza ovunque. Le due non sono mai confuse: trattare un valore ricostruito come se fosse osservato travisarebbe la base evidenziale dell'analisi, ed è evitato per disegno.

\subsection*{Campionamento in profondità}
La griglia verticale RAPID è quasi uniforme (spaziatura 19,33--19,87~m, rapporto max/min di 1,03), non più densa nell'oceano superiore. Nessuna ponderazione per lo spessore degli strati è quindi applicata nella decomposizione: con livelli così uniformemente spaziati, la ponderazione scalerebbe ogni livello di una costante quasi identica, senza effetto sui loading o sulle quote di varianza. Questa è una scelta deliberata e documentata anziché una svista.

% -------------------------------------------------------------------------
\section{Tradurre --- fase di progettazione}
\label{app:limits-translate}

\subsection*{Il vincolo di fedeltà sul design}
Ogni valore nell'artefatto deriva dal bundle analitico, mai da un ridimensionamento inventato. Questo mantiene la rappresentazione fedele anziché illustrativa, a un costo deliberato: l'estetica è vincolata a ciò che i dati sostengono, e dove una scelta artistica confliggerebbe con i dati, prevalgono i dati.

\subsection*{Trasferibilità della grammatica di design}
La grammatica di design mappa proprietà del sistema su canali percettivi attraverso principi cognitivi documentati. La sua trasferibilità è un'affermazione, non una prova: i principi sono generali, ma la loro applicazione fedele dipende dalla specifica struttura che un sistema espone. Una mappatura fedele per la struttura ordinata e a bassa dimensionalità dell'AMOC potrebbe non trasferirsi invariata a un sistema con struttura diversa. La grammatica è offerta come punto di partenza riutilizzabile soggetto ad adattamento al dominio, non come un modello fisso.

% -------------------------------------------------------------------------
\section{Valutare --- fase empirica}
\label{app:limits-evaluate}

\subsection*{Disegno del confronto}
La valutazione confronta due gruppi di partecipanti non specialisti: un gruppo di controllo esposto alla condizione statica convenzionale (C1) e un gruppo sperimentale esposto alle condizioni dinamica ed esperienziale (C2 e C3). Il confronto è quindi tra due intere esperienze di rappresentazione anziché una progressione intra-soggetto attraverso tutte e tre le condizioni. Questa scelta rimuove il confondimento di ordine ed esposizione ripetuta che un disegno sequenziale C1$\rightarrow$C2$\rightarrow$C3 introdurrebbe tra i gruppi, e allinea lo studio direttamente alla domanda di ricerca centrale, che contrappone i formati statici convenzionali a quelli dinamici ed esperienziali.

\subsection*{Forma contro quantità di esposizione}
Poiché il gruppo sperimentale esperisce C2 e C3 mentre il gruppo di controllo esperisce solo C1, i gruppi differiscono nell'estensione dell'esposizione oltre che nel tipo, e il disegno non isola l'effetto della forma esperienziale dall'effetto di un'esposizione più ricca. Confronta due esperienze di rappresentazione nel loro insieme, che è esattamente ciò che la domanda di ricerca chiede. In modo correlato, il blocco sperimentale non isola il contributo dell'interattività (C2) da quello dell'esperienzialità guidata dal machine learning (C3); un minore effetto d'ordine entro il blocco (C2 precede C3) è riconosciuto e non incide sul confronto tra gruppi. Separare questi contributi è lasciato a lavori futuri. Un'ulteriore proprietà del disegno appartiene qui per la stessa ragione: tutti i partecipanti hanno incontrato le tre scene nello stesso ordine fisso, perciò la posizione nella sequenza è confusa con l'identità della scena e ogni confronto \emph{tra} scene porta un possibile contributo da apprendimento o affaticamento. Queste sono proprietà del disegno, riconosciute anziché corrette.

\subsection*{Erogazione asincrona e autoguidata}
Lo studio è erogato come questionario online autoguidato, senza il ricercatore presente. Questa scelta si addice a una valutazione di non specialisti nel proprio ambiente, ma plasma ciò che l'evidenza può essere: cattura il modello mentale che un partecipante esprime per iscritto, non il processo di ragionamento mentre si svolge momento per momento. L'evidenza è quindi l'esito articolato anziché il processo dal vivo, ed è interpretata come tale.

\subsection*{Strumento bilingue}
Il questionario è somministrato in italiano e in inglese, così che i partecipanti rispondano nella propria prima lingua. Le risposte aperte sono analizzate nella lingua in cui sono scritte, anziché tradotte prima della codifica, così che la traduzione non distorca il significato. L'analisi copre quindi due lingue, e i codici sono riconciliati tra esse in fase di reporting.

\subsection*{Peso evidenziale e dimensione campionaria}
Lo studio è esplorativo. Il suo peso evidenziale poggia sull'analisi qualitativa di come i partecipanti articolano la propria comprensione nelle risposte scritte aperte; eventuali item chiusi di comprensione sono riportati in modo descrittivo, non come test inferenziali, in coerenza con la dimensione campionaria realizzata (venticinque partecipanti: dodici nella condizione di controllo e tredici nella sperimentale) e con l'inquadramento esplorativo. Lo studio non rivendica un'inferenza statistica a una popolazione più ampia.

\subsection*{Interpretazione del confronto a due gruppi}
Le aspettative che guidano la fase empirica (E3.1, E3.2) sono formulate attorno al contrasto tra gruppi e non presuppongono la direzione o l'entità di alcun effetto. Differenze qualitative nel modo in cui i due gruppi articolano la propria comprensione possono comparire con o senza differenze accompagnatorie negli item descrittivi di comprensione; entrambi gli esiti sono riportati. Lo studio non è progettato per decidere se tali differenze persisterebbero su scala di popolazione.

\subsection*{La sequenza conclusiva non valutata}
Dopo la parte valutata del questionario, a ogni partecipante sono offerti link a tutte e tre le scene in ogni rappresentazione, per interesse. Le risposte a essa non sono valutate e non entrano nel confronto. Funziona come debriefing e come restituzione al partecipante --- così che chi ha visto solo la versione statica non resti senza l'esperienza --- ed è collocata per disegno fuori dal percorso valutato, così da non poter agire da confondente.

% -------------------------------------------------------------------------
\section{Condotta etica dello studio}
\label{app:ethics-conduct}

La partecipazione è stata informata e volontaria. Ogni partecipante ha ricevuto una scheda informativa in linguaggio piano che spiegava cosa comportava lo studio, quali dati venivano raccolti e come venivano usati, e ha dato il consenso prima di prendere parte; poteva interrompere in qualsiasi momento senza dare una ragione e senza conseguenze. Non è stato usato alcun inganno. Il tema non era sensibile --- i partecipanti ragionavano su una corrente oceanica, non su qualcosa di personale o angosciante --- perciò i principali doveri etici erano quelli standard di consenso, anonimato e protezione dei dati. Le risposte sono state anonimizzate, con i dettagli identificativi rimossi o sostituiti da codici, così che nessun partecipante sia identificabile nella tesi o in alcun output. I dati sono stati conservati in modo sicuro e in linea con il Regolamento generale sulla protezione dei dati, mantenuti solo per il tempo che la ricerca richiedeva e accessibili solo al ricercatore e, dove necessario, al relatore. Lo studio è proceduto solo una volta che il relatore ha confermato che l'approvazione o l'esenzione etica appropriata per uno studio di questa scala era in essere. La scheda informativa e il testo di consenso sono riprodotti nell'Appendice~\ref{app:protocols}.
% -------------------------------------------------------------------------
\section{Uso di assistenti di IA}
\label{app:ai-use}

Nell'interesse della trasparenza, voglio essere chiaro sul ruolo che gli assistenti di intelligenza artificiale hanno avuto nella produzione di questa tesi, e sui confini che vi ho posto attorno.

Gli strumenti di IA sono stati usati come ausilio lungo tutto il lavoro. Claude di Anthropic, incluso l'ambiente Claude Code, mi ha supportato nella stesura e revisione del testo, nella strutturazione delle argomentazioni, nella scrittura e nel debug del codice analitico, e come interlocutore quando avevo bisogno di ragionare su una scelta metodologica. Perplexity è stato usato per individuare letteratura candidata, che ho poi verificato e letto di prima mano prima di citarla. In ogni caso gli strumenti hanno assistito una decisione; non l'hanno presa. Il disegno della ricerca, l'interpretazione dei risultati, le scelte metodologiche e la formulazione finale sono stati miei, e ho riesaminato e deciso tutto ciò che gli strumenti proponevano anziché accettarlo in blocco. Dove l'analisi dipendeva dal giudizio --- la codifica delle risposte dei partecipanti in particolare --- ho trattato ogni suggerimento automatico come un primo passaggio da verificare, non come una risposta. Concretamente: al modello è stata data la griglia di codifica dell'Appendice~\ref{sec:app-coding-grid} e le sue definizioni di categoria, e ha restituito un'etichetta proposta per risposta; ho riesaminato ogni etichetta rispetto alla risposta originale e l'ho corretta ovunque la mia lettura differisse. Il rischio che questa procedura comporta --- che un revisore che vede prima un'etichetta proposta possa ancorarvisi --- è registrato tra i limiti nella \S\ref{sec:eval-limitations}.

Ammetto una piccola ambivalenza, e preferisco dichiararla piuttosto che nasconderla. Una parte di me preferisce il flusso splendidamente imperfetto della scrittura umana, la ruvidezza che porta dentro di sé una persona, e sono consapevole che uno strumento che leviga il linguaggio può anche appiattirlo. Ma usare bene questi strumenti è stato esso stesso una sfida e un esperimento, adatto a una tesi su come la tecnologia possa estendere la comprensione umana anziché sostituirla. L'ho affrontato in quello spirito: come un tentativo di lasciare che lo strumento potenziasse il lavoro mantenendo miei il pensiero, la responsabilità e la voce. Sono soddisfatto di come quell'equilibrio è venuto, e ho cercato di esserne onesto qui anziché silenzioso.
